\begin{abstract}

\begin{center}
{\LARGE Lightweight ResUMamba-based On-Device Real-Time\\
Arrhythmia Classification for Wearable ECG Monitoring}
\end{center}

\begin{flushright}
\parbox[t]{0.9\textwidth}
{
\begin{flushright}
{\bf \txtAuthor}\\
Department of Computer Science\\
The Graduate School, Yonsei University, Korea
\end{flushright}
}
\end{flushright}

\vspace{1em}

Wearable ECG patches enable long-term monitoring of intermittent arrhythmias
during everyday life, but long recordings are difficult to review manually,
and cloud-based analysis introduces additional communication and data-processing
steps. This thesis deploys the existing patient-conditioned ResUMamba beat
classifier on a smartphone and evaluates a lightweight model variant together
with a CPU-oriented deployment approach. The ResU path is modified by replacing
standard convolutions with depthwise separable convolutions and replacing the
decoder's transposed convolutions with parameter-free linear interpolation,
yielding ResDSUMamba, while the Mamba selective scan is rewritten using standard
PyTorch ATen operators and represented as a sequential graph unrolled over the
fixed input length of $T=501$. No new classifier architecture is introduced.

ResDSUMamba has 1.3225 M parameters and 331.38 M MACs, compared with 2.4883 M
and 721.50 M for ResUMamba, corresponding to reductions of 46.9\% and 54.1\%,
respectively.
On a record-wise split of the HiCardi dataset, the Macro-$F_1$ and Micro-$F_1$
of ResDSUMamba are 44.15\% and 83.29\%, compared with 42.28\% and 79.27\% for
ResUMamba, while Macro-Se decreases from 78.00\% to 75.72\%.
On a Galaxy S21+ (Exynos 2100) using a single Cortex-X1 thread in FP32, XNNPACK
delegation reduces the mean $T_{\mathrm{model}}$ of ResUMamba from 1073.57 ms to
48.25 ms, and DSConv reduces it further to 36.70 ms for ResDSUMamba, with a P95
of 38.90 ms.
In a live patch-connected session lasting 309.21 s, 413 beats were processed
at a reported mean heart rate of 80.2 bpm without backlog.
The mean $T_{\mathrm{beat}}$ was 59.4 ms with a P95 of 70.3 ms, while
inference-thread CPU utilization was 7.9\% and peak application-level native
heap usage was 99.2 MB. The live session had no reference labels, so
on-device classification accuracy was not evaluated.

\vspace{\stretch{1}} \noindent
\hrulefill\\
{\bf Key words : }
\parbox[t]{0.8\textwidth}
{Electrocardiogram (ECG), Arrhythmia Classification, State-Space Models (SSMs), On-Device Inference, Mobile CPU Optimization, Wearable Healthcare}

\end{abstract}
\clearpage
